\section{Conclusion}
\label{sec:conclusion}

Reconstructing the uniform-price demand from primitive consumer choices changes the equilibrium domain in the inherited-history Hotelling benchmark of \citet{GehrigShyStenbacka2012}. The accepted Eq.~(12) price vector is the correct smooth switching-branch candidate, but it is a pure-strategy Nash equilibrium only when
\[
x_0\ge
\frac12-\frac{s}{3}
+\frac{\sqrt{3s(s+6)}}6.
\]
Below this boundary no pure-strategy price equilibrium exists; mixed pricing remains outside the scope of the reassessment.

The correction materially narrows the set of parameters over which weak HBP and uniform pricing can be compared as pure equilibria. Direct primitive integration also yields a separate consumer-surplus correction: accepted Eq.~(15) is correct on its switching branch, but Eq.~(16) reverses its sign and the no-switch region requires a different realized-allocation expression. The weak-dominance sign reversal is therefore not reproduced. On the corrected weak pure overlap, HBP raises consumer surplus, lowers the smaller firm's profit, and lowers social welfare relative to uniform pricing. Under strong dominance, consumer-surplus conclusions require an explicit source/nonnegative-margin price selection if below-cost poaching prices are allowed, whereas the welfare comparison is invariant to that selection.

The broader lesson is specific but useful. In price games with history-dependent kinks, an interior candidate should not be labeled an equilibrium until finite deviations across allocation regimes have been ruled out. Here that global check produces an exact threshold, changes the feasible welfare-comparison domain, and separates source conclusions that survive from those that do not.
